% 天文学（综述）
% license CCBYSA3
% type Wiki

本文根据 CC-BY-SA 协议转载翻译自维基百科\href{https://en.wikipedia.org/wiki/Astronomy?utm_source=chatgpt.com}{相关文章}。

\begin{figure}[ht]
\centering
\includegraphics[width=6cm]{./figures/ddcd5c7b339e766a.png}
\caption{欧洲南方天文台（ESO）的帕瑞纳尔天文台（Paranal Observatory）向银河系中心发射激光导星（Laser Guide Star）。} \label{fig_Astron_1}
\end{figure}
天文学是一门研究天体及宇宙中发生的各种现象的自然科学。它运用数学、物理学和化学来解释这些天体的起源及其整体演化。研究对象包括行星、卫星、恒星、星云、星系、流星体、小行星和彗星等；相关的现象则包括超新星爆发、伽马射线暴、类星体、耀变体、脉冲星以及宇宙微波背景辐射等。更广义地说，天文学研究的一切事物都起源于地球大气层之外。而宇宙学（Cosmology）则是天文学的一个分支，专门研究宇宙整体的起源、结构与演化。

天文学是最古老的自然科学之一。在有文字记载的早期文明中，人类就已开始系统地观测夜空。其中包括古埃及人、巴比伦人、希腊人、印度人、中国人、玛雅人，以及美洲的许多原住民。在古代，天文学的范畴十分广泛，涵盖了天体测量学、天体导航、观测天文学以及历法制定等多个领域。

现代专业天文学大体上分为观测天文学与理论天文学两个分支。  
观测天文学的核心任务是通过望远镜与其他仪器对天体进行观测，从而获取数据；然后利用物理学的基本原理对这些数据进行分析。  
理论天文学则侧重于建立计算机模型或解析模型，用以描述天体及其物理现象。  
这两个分支相辅相成：理论天文学旨在解释观测结果，而观测天文学则通过实验数据来验证理论预测。

天文学是少数几个业余爱好者仍然能够发挥积极作用的科学领域之一。这一点在瞬变天体事件（如新星爆发、彗星出现等）的发现与观测中尤为明显。许多重要的天文发现——包括新的彗星——都是由业余天文学家首先发现的。
\subsection{词源}
“天文学”（Astronomy）一词源自希腊语ἀστρονομία（astronomia），由 ἄστρον（astron，“星星”）与 -νομία（-nomia，来自νόμος nomos，“法则”或“规律”）组成，意为“研究天体的学科”。[1]天文学不应与“占星学”（Astrology）混淆——后者是一种信仰体系，声称人类事务与天体位置之间存在关联。二者在起源上有共同渊源，但后来分道扬镳：天文学建立在物理学的科学基础之上，而占星学则缺乏物理学支持。[2]
\subsubsection{“天文学”与“天体物理学”的用法}
在现代用法中，“Astronomy”（天文学）与“Astrophysics”（天体物理学）基本上是同义词。[3][4][5]按照词典定义，天文学是“研究地球大气层之外的物体与物质及其物理和化学性质的科学”；[6]而天体物理学是天文学的一个分支，专门研究“天体的行为、物理性质及其动力过程”。[7]在某些情况下（例如 Frank Shu 所著的《The Physical Universe》（《物理宇宙》）导论中），“天文学”被用来指代该学科的定性研究，而“天体物理学”则强调以物理为导向的定量研究。[8]某些领域（如天体测量学Astrometry）在这种意义上被视为“纯粹的天文学”，而非天体物理学。研究机构或大学的命名方式往往取决于该部门是否历史上隶属于物理学系。[4]事实上，许多职业天文学家持有的是物理学学位，而非天文学学位。[5] 因此，在当代学术语境中，“天文学”与“天体物理学”常被互换使用。[3]
\subsection{历史}
\subsubsection{史前时期}
\begin{figure}[ht]
\centering
\includegraphics[width=6cm]{./figures/2ec12ec85a389120.png}
\caption{尼布拉星象盘（约公元前 1800–1600 年），在一处可能具有天文功能的遗址附近被发现。  该星盘极有可能描绘了太阳或满月、弯月、昴星团（Pleiades）以及夏至与冬至，后两者以盘侧的金条表示。[9][10] 星盘的上缘代表地平线与北方。[11]} \label{fig_Astron_2}
\end{figure}
天文学最初的发展是由实际需求所推动的，例如用于农业活动的历法制定。在有文字记载之前，考古遗址如巨石阵（Stonehenge）就已显示出人类对天文观测的浓厚兴趣。[12]: 15   另一类证据来自文物，例如尼布拉星象盘（Nebra Sky Disc）。这件青铜器被认为是一种天文历法，其设计将一年定义为十二个太阴月（共 354 天），并通过闰月来校正与太阳年的差异。星象盘上镶嵌有象征性的图案，被解释为太阳、月亮和星辰，其中包含一个由七颗星组成的星团。[9][13][14]
\subsubsection{古典时期}
像埃及、两河流域、希腊、印度和中国等文明，在\textbf{跨文化交流}的影响下，纷纷建立了天文观测台，并发展出关于\textbf{宇宙本质}的思想，以及用于测时的\textbf{历法与天文仪器}。[16]天文学早期的一个关键发展是\textbf{巴比伦人}开创了\textbf{数学与科学天文学}，  
为其他文明的天文传统奠定了基础。[17]  
巴比伦人发现了\textbf{月食}以\textbf{沙罗周期（saros cycle）}为规律重复出现，该周期为 223 个朔望月。[18]  

继巴比伦人之后，古希腊与希腊化世界在天文学上取得了重大进展。  
\textbf{希腊天文学}致力于对天体现象寻找\textbf{理性与物理的解释}。[19]  
在公元前 3 世纪，\textbf{萨摩斯的阿里斯塔克斯（Aristarchus of Samos）}估算了月球与太阳的大小与距离，  
并提出了一个以太阳为中心、地球与行星环绕其旋转的\textbf{太阳中心模型（heliocentric model）}。[20]  

在公元前 2 世纪，\textbf{喜帕恰斯（Hipparchus）}计算了月球的大小与距离，  
并发明了已知最早的天文仪器之一，如\textbf{星盘（astrolabe）}。[21]  
他还观测到春分点与秋分点、夏至与冬至相对于恒星位置的微小漂移，  
这一现象即如今所知的\textbf{岁差（precession）}。[12]  
喜帕恰斯还编制了包含\textbf{1020 颗恒星}的星表，  
北半球大多数星座的命名均源自\textbf{希腊天文学}。[22]  

约公元前 150–80 年的\textbf{安提基特拉机械（Antikythera mechanism）}是一种早期的\textbf{模拟计算机}，  
用于计算太阳、月亮与行星在特定日期的位置。  
在其之后，直到 14 世纪欧洲出现机械天文钟，才再次出现与之相当复杂的技术装置。[23]  

在古希腊古典时期之后，天文学长期受\textbf{地心宇宙模型（geocentric model）}或\textbf{托勒密体系（Ptolemaic system）}主导，  
该体系以\textbf{克劳狄乌斯·托勒密（Claudius Ptolemy）}命名。  
他以 13 卷本的著作《\textbf{天文学大成（Almagest）}》（阿拉伯译名）总结了天文学知识，  
并在此后一千多年中成为主要的参考文献。[24]: 196   

在该体系中，地球被认为是\textbf{宇宙的中心}，  
太阳、月亮与恒星围绕地球旋转。[25]  
虽然该理论最终被推翻，但在当时它提供了\textbf{最精确的天体位置预测模型}。[24]
```
